\section{Security}
\label{sec:security}

The security of a mutual-credit system is the question of how much an adversary can make
the community underwrite. Here the question has a closed-form answer, because capacity is
a cut and cuts compose.

Throughout, $S$ denotes a set of accounts and $\partial S$ the set of stakes running
\emph{into} $S$ from outside it. All statements are about a single \emph{membership} ---
one account set, one total order, one underwriter set, one cut. \S\ref{sec:model} explains
why nothing is ever computed across two orders, and \S\ref{sec:oneledger} why there is one.

\subsection{What cannot be manufactured}

Every result in this section applies one statement repeatedly.

\begin{assumption}[The free-signature bound]
\label{ass:free-signatures}
Every ledger event is a signature, and signatures are free. No quantity computed from
ledger events alone can therefore seed insured credit: the seed must be a commitment by a
party with something to lose outside the ledger.
\end{assumption}

Identities are free, settled volume is free, communities are free --- a genesis file costs
nothing --- and a levy on a free quantity is free. The only unforgeable object in the
system is a declared liability, and capacity is the arithmetic of tracing how far one
reaches. Measured against the independent max-flow: a mutual pool funded by a levy on
every settlement, started with no underwriters, accumulates \emph{zero} over $10{,}000$
settled trades, because nothing may be conferred, so nothing is staked, so the levy is a
percentage of nothing. Seeded, the pool does accumulate real stake --- and the members
measured \emph{together with} it still owe exactly the seed: a pool funded from inside
re-labels the seed instead of adding to it, which is why the ledger has no loss pool
(\S\ref{sec:no-pool}). A community that underwrites nothing at genesis can trade
indefinitely on uninsured credit and can never leave that state, since a supply is seated
only by a ceremony and a ceremony needs a seed to be a fraction of.

\subsection{The cut bound}

\begin{theorem}[Cut bound]
\label{thm:cut}
Let $\drawn(S)$ be the total insured credit outstanding to the members of $S$ that is
sourced by underwriters outside $S$. Then
\[
  \drawn(S) \;\le\; \Capa^{0}(S) \;\le\; \sum_{(a,b)\in\partial S} \stake(a,b)
  \;+\; \textstyle\sum_{u\in\uw\setminus S}\supply(u).
\]
\end{theorem}

\begin{proof}
By Definition~\ref{def:reservation} every insured obligation to a member of $S$ holds an
augmenting flow into $S$, and distinct obligations hold disjoint flow, since reservation
subtracts from $\free$ before the next augmentation is computed. The total held is
therefore a feasible flow into $S$ over the gross graph, bounded by the maximum such flow,
which is $\Capa^{0}(S)$. The second inequality is max-flow/min-cut: $\partial S$ together
with the supply arcs of underwriters outside $S$ forms a cut, so the maximum flow is at
most its capacity.
\end{proof}

\begin{remark}[Why the gross reading, and why flow sourced outside $S$]
\label{rem:cut-gross}
Both qualifications are load-bearing, and each fixes a different way the statement is
false without it. Against $\Capa(S)$ the theorem fails for any drawn debtor: the residual
reading nets out exactly the credit the claim is about, so a member who has borrowed to
their ceiling has $\Capa(S)=0$ and $\drawn(S)>0$. And $\drawn$ counts flow sourced
outside $S$ because $\Capa^{0}$ does: an underwriter inside $S$ supplies $S$ nothing in
the measurement, so credit it insures for its own set is neither bounded by this nor
claimed to be. This is what the implementation checks after every committed block
(\S\ref{sec:verification}, invariant~1).
\end{remark}

The proof uses \emph{disjoint flow}, which binds the implementation as much as the
argument: an augmenting path runs from a supply arc through stake edges, and a reservation
must charge both (Remark~\ref{rem:charge-supply}). Charging only the edges leaves the
theorem true of $\Capa$ while the system reaches states it excludes, a bound argued rather
than enforced.

Every subsequent result is a corollary of Theorem~\ref{thm:cut} applied to a well-chosen
$S$. That is the whole of the analysis: a bound that must be re-argued for each new attack
is a bound that will eventually be re-argued wrongly.

\subsection{Identities are free and worthless}

\begin{corollary}[Sybil invariance]
\label{cor:sybil}
Let $x$ be an account with no incident stakes. Then $\Capa(S\cup\{x\}) = \Capa(S)$.
In particular, an account created by an unheralded key has capacity zero, and creating
$N$ of them adds nothing to any coalition.
\end{corollary}

\begin{proof}
$\partial(S\cup\{x\}) = \partial S$, since $x$ contributes no stakes to the boundary,
and $x\notin\uw$ since a supply is seated only by a ceremony
(\S\ref{sec:seed-amendments}) and no ceremony has seated $x$. The two sets induce the
same cuts, hence the same maximum flow.
\end{proof}

There is accordingly no admission decision to make and no rate to limit, which is also why
the ledger carries no creation transition to rate-limit: a row is seated by the first
trade naming its key, and what bounds it is the bond that trade already pays
(\S\ref{sec:implementation}). Account creation is permissionless because it is inert.

\subsection{Fabricated trade is worthless}

\begin{corollary}[Wash invariance]
\label{cor:wash}
Any sequence of acceptances and settlements between members of $S$ leaves $\Capa(S)$
unchanged, regardless of the amounts involved.
\end{corollary}

\begin{proof}
Such settlements modify $\stake(c,d)$ only for $c,d\in S$ by Definition~\ref{def:stake}.
Those stakes have both endpoints in $S$, so they cross no cut into $S$ and $\partial S$
is untouched. By Theorem~\ref{thm:cut} the bound is unchanged; and since the flow must
still cross $\partial S$, so is the maximum.
\end{proof}

Corollary~\ref{cor:wash} is why this design contains no wash-trading detector, no
volume-anomaly heuristic and no rate model: fabricated settlement is not detected and not
penalised, it is arithmetically void. Remark~\ref{rem:not-volume} explains why a design
measuring volume instead of stake cannot have this property: settlement demands two
signatures and no delivery, so its volume is free, and a free quantity placed in
$\partial S$ makes the right-hand side of Theorem~\ref{thm:cut} unbounded.

\subsection{Collusion does not multiply}

\begin{corollary}[Collusion bound]
\label{cor:collusion}
Suppose every stake from $S$ to its complement has its other endpoint at a single
account $m$, and $m$ is not an underwriter. Then $\Capa(S) \le \Capa(\{m\})$.
\end{corollary}

\begin{proof}
Every path into $S$ passes through $m$, so $\{m\}$ is a vertex separator. A feasible flow
into $S$ is therefore a feasible flow into $m$, and the maximum of the former is at most
the maximum of the latter.
\end{proof}

A member who backs an adversary passes on exactly what they hold and no more, however
many accounts they distribute it across or however large the stakes they write. The
community's exposure to a coalition is its exposure to the honest members who chose to
back it, which a community can reason about without knowing who is colluding.

\begin{corollary}[A coalition cannot underwrite itself]
\label{cor:self-uw}
Appointing members of $S$ as underwriters does not raise $\Capa(S)$.
\end{corollary}

\begin{proof}
Definition~\ref{def:capacity} draws supply only from $\uw\setminus S$, so an underwriter
inside $S$ contributes no supply arc to $S$'s own measurement. Their stakes on other
members of $S$ are internal and excluded by Corollary~\ref{cor:wash}.
\end{proof}

\begin{remark}[What this corollary does \emph{not} say]
\label{rem:self-uw-limits}
It is a statement about the members of $S$, and the credit a coalition wants is issued to
the accounts $S$ \emph{backs}, which are outside it and to which $S$'s supply arcs
therefore do count. That is not a gap in the corollary but a question it does not answer,
and the answer comes from elsewhere: a supply is seated only by a ceremony
(\S\ref{sec:seed-amendments}), so $S$ cannot appoint its own accomplices in the first
place. A rule that let a member declare against their own capacity would leave this
corollary true and the community holding $2{,}048$ times its seed in ledger-labelled
insured credit (\S\ref{sec:standing}).
\end{remark}

This is what lets the underwriter role be open to every member without opening a door: a
coalition may propose as many of its own accounts as it likes, measured as a coalition it
gains nothing even if the community endorses them, and measured individually its members
are still bounded by what flows in from outside.

\begin{corollary}[Bounded by declared underwriting]
\label{cor:uw-bound}
For any set $S$, $\Capa(S) \le \sum_{u\in\uw\setminus S}\supply(u)$. In particular an
underwriter's total exposure --- across every member they back, however many --- is at
most their own declared supply.
\end{corollary}

\begin{proof}
The supply arcs of $\uw\setminus S$ form a cut.
\end{proof}

Corollary~\ref{cor:uw-bound} is what makes the recursion in Definition~\ref{def:stake}
terminate, and it is the community's aggregate solvency statement: total insured credit
never exceeds what its underwriters have collectively declared. That total is not a
genesis constant --- it rises when the community endorses an underwriter and falls when
one withdraws, subject to the floor of \S\ref{sec:standing} --- and it is a measure rather
than an upper bound precisely because a supply cannot be declared against standing the
community itself conferred: under a rule that allowed one, \S\ref{sec:standing} measures
the divergence at $204{,}900$ declared against an aggregate of $100$. What the community
can pay is what its underwriters can pay from outside the set being measured.

\subsection{The cost of attack}

\begin{corollary}[Linearity]
\label{cor:linear}
An adversary who acquires genuine backing worth $B$ --- by trading honestly, at honest
prices, with members willing to stake on them --- obtains at most $B$ of capacity across
every account they control.
\end{corollary}

\begin{proof}
Take $S$ to be the adversary's accounts. Their genuine backing is exactly $\partial S$;
apply Theorem~\ref{thm:cut}.
\end{proof}

The proof rests on one word, \emph{genuine}, and that word binds the transition function:
it holds only because \textbf{no transition writes an edge of $\partial S$ without the
signature of somebody outside $S$}. By Remark~\ref{rem:only-settlement} the sole writer is
settlement, which takes the creditor's signature, and a creditor outside $S$ signing is
the adversary having persuaded an outsider to back them, which is what the corollary
counts. That condition is easy to break by accident: a debtor swap looks like a discharge,
and a rule that staked it would let a ring of $k$ members passing one accepted claim
collect an edge each from a creditor who signed once --- $\sum \stake$ linear in $k$ from
one outsider's signature. The swap therefore stakes nothing (\S\ref{sec:medium}), and any
transition that touches the stake graph has to answer the same question.

\begin{corollary}[Seats are linear in backing]
\label{cor:seats}
Let $S'$ be any set of accounts containing no underwriter, and let $\seats(S')$ count the
rows seated by a sponsor in $S'$. Then
\[
  \kConstBondUnitFraction \cdot \vbase \cdot \seats(S')
  \;\le\; \sum_{e \in \partial S'} \peak(e),
\]
where $\peak(e)$ is the all-time maximum weight of arc $e$ and a supply arc counts at its
all-time maximum declaration. Minimising over supersets gives the bound at the cut.
\end{corollary}

\begin{proof}
Every seat holds a flow from the source to its sponsor, of one bond unit, on a reservation
map the credit layer never reads (Property~\ref{prop:seat}). Consider the net seat flow
across $\partial S'$: for a sponsor inside $S'$ the whole amount crosses inward, for one
outside it nets to zero, so the net inward seat flow is exactly
$\kConstBondUnitFraction \cdot \vbase \cdot \seats(S')$. That flow is at most
$\sum_{e \in \partial S'} \seatres(e)$, where $\seatres$ is the seat reservation on an arc,
and each such value is at most $\peak(e)$: a reservation is refused where it would exceed
the arc's weight at the time it is taken, and nothing but a reservation ever raises one.
\end{proof}

Counting by sponsor rather than by row is what the statement needs and the honest reading:
a row an outsider chose to bring in cost that outsider's reach and is not the set's doing,
and a row the set seated is, wherever its account then sits. The bound is over all-time
peaks rather than present weights, and that is exact rather than loose for a permanent
stock backed when it is seated: decay lowers an arc without releasing what a seat took on
it, so renewal to an old peak opens nothing. Time appears nowhere in the statement, which
is the point --- a bound that a wait relaxes is a bound on a rate, and a row is a stock.
Charging the supply arc is what makes the aggregate hold: ten operators each vouched at
$500$ by one underwriter who declared $2{,}500$ hold $125$ seats between them and not
$250$, because the ten stake arcs sum to $5{,}000$ and the arc they all draw through does
not.

There is therefore no leverage anywhere in the system: no configuration of identities,
volumes, timing or collusion converts $B$ of real economic commitment into more than $B$
of underwritten credit. The protocol has no attack surface distinct from its participation
surface. An adversary who wants more must first persuade more members to back them, which
is indistinguishable from becoming a member in good standing, and costs the same.

\subsection{The scope, and how it is honoured}

Everything above is scoped to one ledger: a cut drawn over one order's accounts never
measures a coalition that stood up a second. The scope is narrower than it first appears,
because a ledger is not a community --- two co-operatives sharing an order are inside one
cut, and every result above applies to a coalition spanning both --- and it is honoured by
communities that expect to trade sharing a ledger. What lies outside the order lies outside
the arithmetic; the alternatives were measured (\S\ref{sec:oneledger}) and each voids the
corollaries above at a stroke.

\subsection{What is assumed}

Theorem~\ref{thm:cut} and its corollaries assume no honesty on the part of any account, no
bound on identity creation, no external registry, and no observation of behaviour. They
assume:

\begin{enumerate}[leftmargin=1.6em]
  \item \textbf{Underwriters can bear what they declared.} A supply is a promise; the
        protocol enforces its consequences (\S\ref{sec:recourse}) but cannot verify in
        advance that an underwriter is good for it. A community that lets members declare
        more than they can bear has mispriced itself, and no later mechanism corrects
        that.
  \item \textbf{The order is agreed.} Reservation is a conserved budget, so replicas that
        disagree on the sequence of acceptances disagree on who holds it
        (Remark~\ref{rem:order}).
  \item \textbf{Capacity is computed identically everywhere.} The capacity path is
        integer throughout for this reason; see \S\ref{sec:implementation}.
  \item \textbf{No quantity from another order enters.} A single foreign input --- a
        recognised supply, a discharge authorised by a foreign root --- voids
        Corollaries~\ref{cor:sybil}, \ref{cor:wash}, \ref{cor:self-uw}
        and~\ref{cor:linear} together. The assumption is discharged by there being no
        transition that can carry such an input.
\end{enumerate}

Nothing else is assumed, and in particular nothing about the proportion of honest members,
the cost of an identity, or the quality of any membrane.

\subsection{What this does not defend}

\paragraph{A creditor's own judgement.} Uninsured credit is unbounded by construction
(\S\ref{sec:standing}). A member who lends beyond a debtor's capacity can lose everything
they lent, and the protocol neither prevents it nor compensates it: the decision sits with
the party holding the information.

\paragraph{Backing withdrawn by decay.} A member's capacity can fall without any act of
their own --- because the stakes behind it decay, because those who staked lost their own
capacity, or because an underwriter reduced their supply. Standing is a claim on others'
present willingness, not a possession, and a community whose members all stop renewing
watches its credit fall to zero (Remark~\ref{rem:winddown}).

\paragraph{Allocation under contention.} When several creditors would lend to the same
debtor, capacity is taken first-come-first-served: the total is bounded correctly, but
which creditor obtains it is decided by transaction order rather than by any notion of
fairness, and order within a block is chosen by its proposer. A validator that sees a
pending acceptance which would exhaust an underwriter can place its own ahead of it in the
same block. It is the one place a validator's private interest touches who is insured; what
it can take is bounded by the same cut as everybody else's, and a validator is an
institution the community chose. Ordering a block's acceptances by their digest would not
remove the lever but auction it: a submitter controls the envelope, so anyone can grind a
digest to the front of a block in milliseconds, where today only the proposer reorders;
and ordering by a value no submitter knows in advance hands the choice back to the
proposer, who assembles the batch. The statement above is the honest one.

\paragraph{Acceptance.} No result in this section implies that anybody will accept an
obligation in this medium. What is proved is that it cannot be inflated by an adversary,
which is necessary and not sufficient; the motive to accept comes from the clearing itself
(\S\ref{sec:sale}) and is bounded by the condition under which clearing is possible at all
(\S\ref{sec:adoption}), and neither is a theorem about an adversary. Nor does any result
supply a newcomer with their first willing counterparty (Remark~\ref{rem:friction}).
